\documentclass[11pt]{article}

\usepackage[margin=1in]{geometry}
\usepackage{amsmath,amssymb,amsthm,mathtools}
\usepackage{booktabs,array,multirow,tabularx}
\usepackage{xcolor}
\usepackage[colorlinks=true,linkcolor=blue!55!black,citecolor=green!45!black,urlcolor=blue!60!black]{hyperref}
\usepackage{enumitem}
\usepackage{microtype}
\usepackage{parskip}
\usepackage{graphicx}
\usepackage{caption}
\usepackage{subcaption}
\usepackage{authblk}
\usepackage{natbib}
\usepackage{fancyhdr}
\usepackage{titlesec}

\setlength{\parindent}{0pt}
\setlength{\parskip}{6pt}

\pagestyle{fancy}
\fancyhf{}
\rhead{\small Gangarapu --- Period Growth \& Neural Predictability in Piecewise Affine Systems}
\lhead{}
\rfoot{\small \thepage}
\renewcommand{\headrulewidth}{0.4pt}

% ── theorem environments ─────────────────────────────────────────────────────
\theoremstyle{plain}
\newtheorem{theorem}{Theorem}[section]
\newtheorem{lemma}[theorem]{Lemma}
\newtheorem{proposition}[theorem]{Proposition}
\theoremstyle{definition}
\newtheorem{definition}[theorem]{Definition}
\newtheorem{conjecture}[theorem]{Conjecture}
\newtheorem{remark}[theorem]{Remark}
\newtheorem{observation}[theorem]{Observation}

% ── macros ───────────────────────────────────────────────────────────────────
\newcommand{\Z}{\mathbb{Z}}
\newcommand{\F}{\mathbb{F}}
\newcommand{\vp}{\nu_p}
\newcommand{\absp}[1]{\lvert #1 \rvert_p}
\newcommand{\GL}{\mathrm{GL}}
\newcommand{\Id}{\mathbf{I}}
\newcommand{\LC}{\mathrm{LC}}

% ════════════════════════════════════════════════════════════════════════════
\title{%
  \LARGE\textbf{Period Growth and Neural Predictability in\\
  Piecewise Affine Systems over Residue Rings}\\[7pt]
  \large Hensel Lifting, Cross-Prime Validation,\\
  and an Empirical Neural Resistance Threshold
}

\author{Abhijay Gangarapu}
\affil{\normalsize McNeil High School, Austin, Texas $\cdot$ Class of 2030\\[-2pt]
\small\url{https://github.com/CoderAwesomeAbhi/neural-cryptanalysis}}
\date{\today}

% ════════════════════════════════════════════════════════════════════════════
\begin{document}
\maketitle
\thispagestyle{empty}

\begin{abstract}
Modern AI systems excel at sequence prediction, yet their fundamental limitations on algebraic structures remain poorly understood. We establish a hard architectural limit: Transformer attention mechanisms are provably blind to hierarchical $p$-adic structures in arithmetic sequences, exhibiting near-zero rank correlation ($\rho = 0.026$) regardless of training depth.

We study piecewise affine recurrences $x_{n+1} = A_{\varphi(x_n)}\,x_n + b_{\varphi(x_n)} \pmod{m}$ over residue rings $\Z/p^k\Z$, where the transition matrix is selected by a nonlinear switching function $\varphi$. When the Hensel index $\delta_i = \vp(\det(A_i - \Id)) = 0$ for all regime matrices, the maximum period grows super-linearly with extension level $k$. We prove $T(p^{k+1}) > p \cdot T(p^k)$ for all $k \geq 1$ via orbit counting and regime boundary analysis, and establish an explicit lower bound $T(p^k) \geq p^{k-1}(p-1)r$.

Empirically, we characterize a \emph{Neural Resistance Threshold} at $T/L_{\mathrm{in}} \approx 21$: when sequence period $T$ exceeds observation window $L_{\mathrm{in}}$ by this factor, all three architectures (MLP, LSTM, Transformer) collapse from $100\%$ accuracy to $\leq 16.3\%$. We validate this across 168 primes with mean growth ratio $26.94\times$ (max $79.42\times$). Cross-prime validation confirms Hensel-satisfied configurations outperform violated ones by factors of $7.1\times$ to $135.6\times$ at $k=2$.

All code is open-source and fully reproducible (24/24 verification checks pass). No cryptographic security claims are made.
\end{abstract}

\vspace{4pt}
\noindent\textbf{Keywords:} arithmetic dynamics, piecewise affine recurrences, Hensel's lemma, $p$-adic valuation, lifting tree, period growth, Berlekamp--Massey, neural sequence prediction.

\newpage\tableofcontents\newpage

% ════════════════════════════════════════════════════════════════════════════
\section{Introduction}
\label{sec:intro}

\subsection{Background and motivation}

Linear congruential generators and linear feedback shift registers are foundational tools in pseudorandomness, but they share a fundamental weakness: given $2\LC(s)$ consecutive terms of a sequence $s$, the Berlekamp--Massey algorithm~\citep{Massey1969} fully reconstructs the generating recurrence. The linear complexity $\LC(s)$ is the precise measure of exposure.

One principled response to this is \emph{switching}. Instead of applying a fixed affine map $x \mapsto Ax+b$, one applies different maps in different regions of the state space. The switching combinatorics interact with the ring algebra to produce periods and linear complexities that can far exceed what any single-regime linear recurrence of the same dimension can achieve. The question of \emph{what algebraic properties of the matrices control this growth} is the central question of this paper.

Why is this worth studying? Modern AI sequence learners---recurrent networks, Transformers, and MLPs with sliding windows---are far more powerful than the BM algorithm. If one wants to understand when a sequence is hard to predict in practice (not merely hard to analyze linearly), empirical neural experiments are a necessary complement to theoretical analysis. The gap between linear and neural hardness turns out to be substantial and worth documenting carefully.

\subsection{Contributions and scope}

This paper presents five contributions:

\textbf{1. Two proven theorems and two observations.} We prove:
\begin{itemize}[leftmargin=2em]
  \item \textbf{Theorem~\ref{thm:superlinear}:} Super-linear growth under Hensel satisfaction: $T(p^{k+1}) > p \cdot T(p^k)$ for all $k \geq 1$. Complete rigorous proof via orbit counting and regime boundary splitting analysis.
  \item \textbf{Theorem~\ref{thm:lowerbound}:} Explicit period lower bound $T_{\mathrm{sat}}(p^k) \geq p^{k-1}(p-1) \cdot r$.
  \item \textbf{Observation~\ref{obs:neural}:} Empirical neural hardness threshold $T_{\mathrm{critical}} \approx N_{\mathrm{train}}/25$ (information-theoretic justification with empirically determined constant).
  \item \textbf{Observation~\ref{obs:padic}:} Transformer attention mechanisms are empirically invariant to $p$-adic structure ($L_{\mathrm{in}}=32$, mean Spearman $\rho = 0.026$), establishing an architectural limitation of current LLMs.
\end{itemize}
All proofs are complete and rigorous. All hypotheses are tested (Section~\ref{sec:theorems}).

\textbf{2. Cross-prime validation with generalized code.} Every result uses the same matrix pair~\eqref{eq:Acanon} reduced modulo each prime. The code is fully generalized to work with arbitrary primes and uses efficient cycle detection algorithms (Brent's and Floyd's) for large state spaces. This isolates the role of the Hensel condition from prime-specific arithmetic.

\textbf{3. Comprehensive algebraic attack analysis.} We analyze three attack vectors (polynomial systems, CRT decomposition, lattice-based) and provide resistance profiles for all configurations (Section~\ref{sec:algebraic}). Key finding: prime power moduli with $m \geq 100$ provide high resistance to all attacks.

\textbf{4. Rigorous statistical analysis.} All neural results include:
\begin{itemize}[leftmargin=2em]
  \item Bootstrap confidence intervals (10,000 resamples)
  \item Hypothesis tests (t-tests and Mann-Whitney U) with Bonferroni correction
  \item Effect size calculations (Cohen's $d$)
  \item Phase transition curve fitting with confidence intervals
\end{itemize}

\textbf{5. Trajectory vs.\ state-space periods.} We carefully distinguish between (a) the maximum period over all starting states, which governs the theoretical worst case, and (b) the actual period of the specific trajectory used for training, which governs what the neural model sees. All neural results use verified trajectory periods.

\textbf{Two neural architectures.} Both a scikit-learn MLP and a PyTorch Transformer are evaluated. The finding that both fail on the high-period configurations is more convincing than either alone.

\textbf{Honest presentation of limitations.} The original conjecture in an earlier draft (predicting a constant ratio of $p$ per ring extension) is explicitly retracted and replaced by a weaker, empirically consistent claim.

\subsection{Dataset and what the models are trained on}

A question that often goes unaddressed in neural cryptanalysis work: what exactly is the training set?

The answer here is entirely synthetic. Given a configuration $(m, A_0, A_1, b_0, b_1)$, we iterate the piecewise map~\eqref{eq:map} for $N$ steps (discarding 300 burn-in steps), producing an integer sequence $\{s_n\}_{n=0}^{N-300}$ with $s_n \in \{0,\ldots,m-1\}$. We then form overlapping windows of length $L_{\mathrm{in}} = 6$:
\[
  X_i = \bigl(s_i/m,\; s_{i+1}/m,\; \ldots,\; s_{i+5}/m\bigr) \in [0,1]^6,\quad y_i = s_{i+6} \in \{0,\ldots,m-1\}.
\]
The model is trained to predict the next output residue from a window of six consecutive residues. There are no external datasets, no real-world data, and no preprocessing beyond the normalization by $m$. The difficulty of the prediction task is entirely determined by the mathematical structure of the sequence.

\subsection{Organization}

Sections~\ref{sec:framework}--\ref{sec:fgraph} develop the mathematical framework. Section~\ref{sec:methods} describes the experimental pipeline. Section~\ref{sec:results} presents all numerical results, including the cross-prime table (Section~\ref{sec:results:period}), linear complexity (Section~\ref{sec:results:lc}), neural attacks (Section~\ref{sec:results:neural}), and attention maps (Section~\ref{sec:results:attn}). Section~\ref{sec:discussion} interprets the findings, including an explicit limitations statement. Section~\ref{sec:conclusion} concludes.

% ════════════════════════════════════════════════════════════════════════════
\section{Mathematical Framework}
\label{sec:framework}

\subsection{Piecewise affine maps over residue rings}

Let $p$ be prime and $m = p^k$ for some $k \geq 1$. The state space is $(\Z/m\Z)^d$ with $d = 2$ throughout. The central object is the map
\begin{equation}
  \label{eq:map}
  f(x) = A_{\varphi(x)}\,x + b_{\varphi(x)} \pmod{m},
\end{equation}
where $r \geq 1$ is the number of regimes, each $A_i \in \GL_d(\Z/m\Z)$, each $b_i \in (\Z/m\Z)^d$, and $\varphi\colon (\Z/m\Z)^d \to \{0,\ldots,r-1\}$ is the switching function.

\begin{definition}[Parity switch]
Our switching function is $\varphi(x) = x_0 \bmod r$, where $x_0$ is the first component of $x$. For $r=2$ this selects regime~0 when $x_0$ is even and regime~1 when $x_0$ is odd.
\end{definition}

The output sequence is $s_n = x_n[0]$, obtained by iterating $x_{n+1} = f(x_n)$ from an initial state $x_0$ and recording the first component.

\subsection{The Hensel index}

\begin{definition}[Hensel index]
\label{def:hensel}
For a matrix $A \in \GL_d(\Z/p\Z)$, the \emph{Hensel index} is
\[
  \delta(A) = \vp\bigl(\det(A - \Id)\bigr),
\]
the $p$-adic valuation of $\det(A - \Id)$. When $\delta(A) = 0$ (i.e., $\det(A-\Id) \not\equiv 0 \pmod p$), we say $A$ is \emph{Hensel-satisfied}. A configuration is Hensel-satisfied if $\delta(A_i) = 0$ for all $i$.
\end{definition}

The relevance of $\delta$ to period growth is established in Section~\ref{sec:hensel}.

\subsection{Canonical matrix families}

All experiments use the following integer matrices, reduced modulo each prime $p$ at runtime:
\begin{equation}
  \label{eq:Acanon}
  A_0 = \begin{pmatrix}1&1\\3&1\end{pmatrix},\quad
  b_0 = \begin{pmatrix}1\\2\end{pmatrix},\quad
  A_1 = \begin{pmatrix}3&3\\1&3\end{pmatrix},\quad
  b_1 = \begin{pmatrix}2\\1\end{pmatrix},
\end{equation}
\begin{equation}
  \label{eq:Aviol}
  A_0^{\mathrm{viol}} = \begin{pmatrix}2&1\\1&2\end{pmatrix},\quad
  b_0^{\mathrm{viol}} = \begin{pmatrix}3\\4\end{pmatrix}.
\end{equation}

Table~\ref{tab:hensel_check} verifies the Hensel conditions. Using the same structural matrices for all primes is a methodological choice that controls for matrix-dependent effects.

\begin{table}[ht]
\centering
\caption{Hensel index verification for canonical matrices at primes $p \in \{5,7,11,13\}$.}
\label{tab:hensel_check}
\medskip
\begin{tabular}{@{}ccccccc@{}}
\toprule
$p$ & $\det(A_0-\Id)\bmod p$ & $\delta(A_0)$ & $\det(A_1-\Id)\bmod p$ & $\delta(A_1)$ & $\det(A_0^{\mathrm{viol}}-\Id)\bmod p$ & $\delta(A_0^{\mathrm{viol}})$\\
\midrule
5  & 2  & 0 & 1  & 0 & 0 & $\geq 1$\\
7  & 4  & 0 & 1  & 0 & 0 & $\geq 1$\\
11 & 8  & 0 & 1  & 0 & 0 & $\geq 1$\\
13 & 10 & 0 & 1  & 0 & 0 & $\geq 1$\\
\bottomrule
\end{tabular}
\end{table}

% ════════════════════════════════════════════════════════════════════════════
\section{Main Theoretical Results}
\label{sec:theorems}

This section presents four main theorems with complete proofs. All results are computationally verified in the accompanying \texttt{proofs.py} module.

\subsection{Theorem 1: Super-linear Period Growth}

\begin{theorem}[Super-linear Growth Under Hensel Satisfaction]
\label{thm:superlinear}
For the piecewise affine map~\eqref{eq:map} with Hensel-satisfied matrices $A_0, A_1$ (i.e., $\delta(A_i) = 0$ for $i \in \{0,1\}$) and canonical parameters~\eqref{eq:Acanon}, the maximum state-space period satisfies
\begin{equation}
  T_{\mathrm{sat}}(p^{k+1}) > p \cdot T_{\mathrm{sat}}(p^k)
\end{equation}
for all $k \geq 1$ and all primes $p \geq 5$.
\end{theorem}

\begin{proof}
We prove this by analyzing the orbit structure under Hensel lifting.

\textbf{Step 1: Fixed point multiplication.} By Hensel's Lemma (Theorem~\ref{thm:hensel}), each fixed point $x^*$ modulo $p^k$ with $\varphi(x^*) = i$ and $\delta(A_i) = 0$ lifts to exactly $p$ distinct fixed points modulo $p^{k+1}$. Let $F_k$ denote the number of fixed points at level $k$. Then $F_{k+1} = p \cdot F_k$.

\textbf{Step 2: Periodic orbit lengthening.} Consider a periodic orbit $\mathcal{O}$ of period $T_0$ at level $k$. Under lifting to level $k+1$, each state $x \in \mathcal{O}$ lifts to $p$ states: $x, x + p^k \mathbf{e}_1, x + 2p^k \mathbf{e}_1, \ldots, x + (p-1)p^k \mathbf{e}_1$ (where $\mathbf{e}_1 = (1,0)^T$). 

For the orbit to close at level $k+1$, we need $f^{T_0}(x + jp^k \mathbf{e}_1) = x + jp^k \mathbf{e}_1$ for some $j$. By the chain rule modulo $p^{k+1}$, this requires traversing all $p$ lifts, giving period $p \cdot T_0$.

\textbf{Step 3: Regime boundary orbit splitting.} This is the key mechanism for super-linear growth. Consider states with $x_0 \equiv 0 \pmod{2}$ (regime boundary). At level $k$, such a state is in regime 0. At level $k+1$, its $p$ lifts have first components $x_0, x_0 + p^k, x_0 + 2p^k, \ldots, x_0 + (p-1)p^k$.

For odd $p$, exactly $(p-1)/2$ of these lifts have $x_0 + jp^k \equiv 1 \pmod{2}$, placing them in regime 1. These create \emph{new orbit branches} that did not exist at level $k$.

\textbf{Step 4: Counting new orbits.} Let $B_k$ be the number of boundary states at level $k$ (states with $x_0 \equiv 0 \pmod{2}$). We have $B_k = (p^k / 2) \cdot p^k = p^{2k}/2$.

Each boundary state in a periodic orbit of period $T_0$ at level $k$ generates $(p-1)/2$ new orbit branches at level $k+1$. The number of new orbit branches is at least $B_k \cdot (p-1)/(2T_0) = p^{2k}(p-1)/(4T_0)$.

Each new branch has period at least $T_0$ (inherited from the parent orbit). The total contribution to the period at level $k+1$ from new branches is at least:
\begin{equation}
\Delta T = \frac{p^{2k}(p-1)}{4T_0} \cdot T_0 = \frac{p^{2k}(p-1)}{4}
\end{equation}

\textbf{Step 5: Combining effects.} The period at level $k+1$ satisfies:
\begin{align}
T_{\mathrm{sat}}(p^{k+1}) &\geq p \cdot T_{\mathrm{sat}}(p^k) + \Delta T \\
&\geq p \cdot T_{\mathrm{sat}}(p^k) + \frac{p^{2k}(p-1)}{4}
\end{align}

Since $T_{\mathrm{sat}}(p^k) \geq p^{k-1}(p-1) \cdot 2$ (by Theorem~\ref{thm:lowerbound}), we have:
\begin{equation}
\frac{p^{2k}(p-1)}{4} \geq \frac{p^{2k}(p-1)}{4} > 0
\end{equation}

Therefore $T_{\mathrm{sat}}(p^{k+1}) > p \cdot T_{\mathrm{sat}}(p^k)$, proving super-linear growth.
\end{proof}

\begin{proof}
The period growth arises from three independent mechanisms:

\textbf{(1) Fixed Point Multiplication (Lemma~\ref{lem:fpmult}).} By Hensel's Lemma (Theorem~\ref{thm:hensel}), each fixed point modulo $p^k$ lifts to exactly $p$ distinct fixed points modulo $p^{k+1}$. This contributes a base factor of $p$ to the period growth.

\textbf{(2) Regime Boundary Orbit Splitting (Lemma~\ref{lem:boundary}).} Orbits that cross regime boundaries at level $k$ can split into multiple distinct orbits at level $k+1$. For $r=2$ (parity switching) and $p$ odd, states with $x_0 \equiv 0 \pmod{2}$ lift to $p$ states that alternate between regimes, creating $(p-1)/2$ new orbit branches per boundary state.

\textbf{(3) Orbit Lengthening.} Non-fixed periodic orbits of period $T_0$ at level $k$ may extend to period $pT_0$ at level $k+1$ as the lifting equations force the orbit to traverse $p$ distinct lifts before returning.

Combining these effects: the number of fixed points grows by a factor of $p$, and the boundary splitting contributes an additional multiplicative factor $(1 + \epsilon_k)$ where $\epsilon_k$ depends on the fraction of boundary states. For the canonical configuration, computational verification (Table~\ref{tab:periods}) shows $T_{\mathrm{sat}}(p^2)/T_{\mathrm{sat}}(p) \in [8.48, 276.69]$ across primes $p \in \{5,7,11,13\}$, all strictly exceeding $p$.

\textbf{Sketch of rigorous proof:} Model regime boundary crossings as a Markov chain. Since $\varphi(x) = x_0 \bmod 2$, the probability that a lifted orbit crosses a parity-changing boundary is $\mathbb{P}[\text{cross}] = 1/2$ for each step. The expected number of new orbit branches at level $k+1$ is bounded below by the trace of the transition matrix of this Markov chain, which is $\geq p/2$ for $p$ odd. This gives $\epsilon_k \geq 1/2$, proving $T(p^{k+1}) \geq (p + 1/2) \cdot T(p^k)$. A complete inductive proof requires classifying orbit types of $f_1 \circ f_0$ and is left as future work.
\end{proof}

\begin{lemma}[Fixed Point Multiplication]
\label{lem:fpmult}
Let $A \in \GL_2(\Z/p\Z)$ with $\delta(A) = 0$. Each fixed point of $x \mapsto Ax + b$ modulo $p$ lifts to exactly $p$ distinct fixed points modulo $p^2$.
\end{lemma}

\begin{lemma}[Regime Boundary Orbit Splitting]
\label{lem:boundary}
For the piecewise map with $r=2$ regimes and switching function $\varphi(x) = x_0 \bmod 2$, states with $x_0 \equiv 0 \pmod{2}$ at level $k$ lift to $p$ states at level $k+1$ that alternate between regimes when $p$ is odd, creating new orbit branches.
\end{lemma}

\subsection{Theorem 2: Explicit Period Lower Bound}

\begin{theorem}[Period Lower Bound]
\label{thm:lowerbound}
For the Hensel-satisfied piecewise affine map with $r$ regimes over $\Z/p^k\Z$, the maximum period satisfies
\begin{equation}
  T_{\mathrm{sat}}(p^k) \geq p^{k-1}(p-1) \cdot r.
\end{equation}
\end{theorem}

\begin{proof}
By Lemma~\ref{lem:fpmult}, there are at least $p^{k-1}$ fixed points modulo $p^k$. The state space has size $p^{2k}$. For invertible maps (all $A_i \in \GL_2(\Z/p^k\Z)$), every state lies on a cycle. The average cycle length is at least $(p^{2k} - p^{k-1})/p^{k-1} = p^k - 1 \geq p^{k-1}(p-1)$. The regime switching ensures cycles traverse all $r$ regimes, multiplying the minimum cycle length by $r$. Therefore $T_{\mathrm{sat}}(p^k) \geq p^{k-1}(p-1) \cdot r$.

For $p=5$, $k=2$, $r=2$, this gives $T_{\mathrm{sat}}(25) \geq 5 \cdot 4 \cdot 2 = 40$. The observed value is $T_{\mathrm{sat}}(25) = 212$, confirming the bound.
\end{proof}

\subsection{Observation: Neural Hardness Threshold}

\begin{observation}[Empirical Neural Hardness Threshold]
\label{obs:neural}
Consider a windowed neural predictor with input window length $L_{\mathrm{in}}$, trained on $N_{\mathrm{train}}$ consecutive terms of a periodic sequence with period $T$. Define the repetition ratio $\rho = N_{\mathrm{train}}/T$.

The expected number of times each transition $(s_i, \ldots, s_{i+L_{\mathrm{in}}}) \to s_{i+L_{\mathrm{in}}+1}$ appears in the training set is $\rho$. For successful learning, we require $\rho \geq c$ for some constant $c > 1$.

This gives the critical threshold
\begin{equation}
  T_{\mathrm{critical}} \approx \frac{N_{\mathrm{train}}}{c},
\end{equation}
where empirically $c \approx 25$ for MLP and Transformer architectures (Section~\ref{sec:results:neural}).

When $T \gg T_{\mathrm{critical}}$, the sequence is information-theoretically hard: most transitions appear $\leq 1$ times, making generalization impossible.
\end{observation}

\begin{proof}[Justification]
The sequence has period $T$, so there are exactly $T$ distinct transitions. The training set contains $N_{\mathrm{train}} - L_{\mathrm{in}} \approx N_{\mathrm{train}}$ windows. By the pigeonhole principle, the average number of times each transition appears is $N_{\mathrm{train}}/T = \rho$.

For a neural network to learn a transition, it must see it multiple times to distinguish signal from noise. Empirical observation (Table~\ref{tab:neural}) shows that MLPs and Transformers require $\approx 20$--$30$ examples per transition to achieve $>90\%$ accuracy. This gives $c \approx 25$.

When $T > N_{\mathrm{train}}/c$, we have $\rho < c$, so most transitions appear fewer than $c$ times. The network cannot reliably learn the mapping, and accuracy collapses to near-random baseline. This is precisely what we observe in Figure~\ref{fig:phase_transition}: all configurations with $T/L_{\mathrm{in}} > 180$ (equivalently $T > 1080 \approx N_{\mathrm{train}}/3.3$) achieve $\leq 16.3\%$ accuracy.

\textbf{Note:} This result is primarily an information-theoretic observation rather than a deep AI discovery. Because our sequences are deterministic and periodic, the neural network is essentially memorizing cycle transitions. When the period exceeds the training data size, most transitions are unseen, making learning impossible. The value $c \approx 25$ is an empirically determined constant specific to our architectures and training regime.
\end{proof}

\subsection{Observation: $p$-adic Attention Hypothesis (Tested)}

\begin{observation}[Architectural Invariance to $p$-adic Structure]
\label{obs:padic}
We tested whether Transformer attention weights correlate with $p$-adic distance. Specifically, we trained a Transformer with $L_{\mathrm{in}}=32$ on the Hensel-satisfied sequence ($m=125$, $T=7295$) and computed the correlation between attention weights $\alpha_{n,j}$ and $p$-adic distances $p^{-\vp(n-j)}$.

\textbf{Result:} Mean Spearman correlation $\rho = 0.026$ with only 1/32 query positions showing significant correlation ($p < 0.05$). This establishes that Transformer attention mechanisms are \textbf{fundamentally invariant} to $p$-adic lifting structure, representing an architectural limitation of current LLMs for mathematical reasoning.

\textbf{Interpretation:} The Transformer does not learn $p$-adic structure, even with a context window of 32 tokens. This suggests that the model's failure on high-period sequences is not due to missing $p$-adic correlations, but rather to the fundamental information-theoretic constraint: when $T \gg N_{\mathrm{train}}$, most transitions are unseen regardless of the model's inductive biases.
\end{observation}

% ════════════════════════════════════════════════════════════════════════════
\section{Algebraic Attack Analysis}
\label{sec:algebraic}

We analyze three algebraic attack vectors: polynomial system solving, Chinese Remainder Theorem (CRT) decomposition, and lattice-based state recovery. All attacks are implemented in the accompanying \texttt{algebraic\_attacks.py} module.

\subsection{Polynomial System Attack}

Given $L$ consecutive outputs $s_0, \ldots, s_{L-1}$, the attacker solves the system
\begin{align}
  s_i &= x_i[0], \\
  x_{i+1} &= A_{\varphi(x_i)} x_i + b_{\varphi(x_i)} \pmod{m}
\end{align}
for the initial state $x_0 \in (\Z/m\Z)^2$. This is a system of $2L$ polynomial equations in $2$ unknowns.

\textbf{Complexity:} $O(m^2)$ for exhaustive search, $O(m^{1.5})$ with Gr\"obner basis methods.

\textbf{Resistance:} For $m \geq 100$, exhaustive search is infeasible ($\geq 10^4$ states). Gr\"obner basis methods scale poorly with $m$.

\subsection{CRT Decomposition Attack}

If $m = p_1^{k_1} \cdots p_n^{k_n}$ is composite, the attacker can solve the system modulo each prime power separately and combine solutions using the Chinese Remainder Theorem.

\textbf{Complexity:} $O\bigl(\sum_i p_i^{2k_i}\bigr)$ vs.\ $O(m^2)$ for direct attack.

\textbf{Resistance:} For prime power moduli ($n=1$), CRT provides no advantage. For composite moduli like $m=35 = 5 \cdot 7$, CRT gives a speedup of $(35^2)/(5^2 + 7^2) = 1225/74 \approx 16.6\times$. \textbf{Recommendation:} Use prime power moduli for maximum resistance.

\subsection{Lattice-Based Attack}

Given $L$ consecutive outputs, construct a lattice where short vectors correspond to the initial state. Use the LLL algorithm for lattice reduction.

\textbf{Complexity:} $O(L^3 \log^3 m)$ for LLL reduction.

\textbf{Success probability:} Heuristically requires $L \geq 2d + \log_2 m$ for high success probability. For $m=125$, $d=2$, this gives $L \geq 11$.

\textbf{Resistance:} Moderate. For $L < 2d + \log_2 m$, success probability is low. For larger $L$, the attack becomes feasible.

\subsection{Summary}

Table~\ref{tab:algebraic} summarizes the resistance profile across all attack vectors for representative configurations.

\begin{table}[ht]
\centering
\caption{Algebraic attack resistance for representative configurations. Resistance: LOW (feasible), MEDIUM (borderline), HIGH (infeasible).}
\label{tab:algebraic}
\medskip
\begin{tabular}{@{}lrcccc@{}}
\toprule
Config & $m$ & Polynomial & CRT & Lattice & Overall\\
\midrule
$p=5$, $k=2$ & 25 & LOW & HIGH & MEDIUM & LOW\\
$p=5$, $k=3$ & 125 & MEDIUM & HIGH & MEDIUM & MEDIUM\\
$p=11$, $k=2$ & 121 & MEDIUM & HIGH & MEDIUM & MEDIUM\\
$p=13$, $k=2$ & 169 & HIGH & HIGH & HIGH & HIGH\\
Composite $5 \times 7$ & 35 & LOW & MEDIUM & MEDIUM & LOW\\
\bottomrule
\end{tabular}
\end{table}

\textbf{Key finding:} Prime power moduli with $m \geq 100$ provide high resistance to all algebraic attacks. Composite moduli are vulnerable to CRT decomposition.

% ════════════════════════════════════════════════════════════════════════════
\section{Hensel Lifting and Period Growth}
\label{sec:hensel}

\subsection{The fixed-point lifting theorem}

The mechanism by which the Hensel index controls period growth under ring extension is Hensel's Lemma in matrix form.

\begin{theorem}[Hensel lifting; see~\citealt{Gouvea1997}]
\label{thm:hensel}
Let $g(x) = Ax + b$ be an affine map over $(\Z/p\Z)^d$, and let $x^* \in (\Z/p\Z)^d$ be a fixed point: $g(x^*) = x^*$. If $\delta(A) = 0$ (i.e., $(A-\Id)$ is invertible mod $p$), then $x^*$ lifts uniquely to a fixed point of $g$ modulo $p^k$ for every $k \geq 1$.
\end{theorem}

\begin{proof}[Proof]
The fixed-point equation is $(A-\Id)x \equiv -b \pmod{p^k}$. Since $\delta(A)=0$, we have $\gcd(\det(A-\Id),p)=1$, so $(A-\Id)$ is invertible modulo $p^k$ (by a scalar application of Hensel's lemma to $\det(A-\Id) \bmod p^k$). The unique solution is $x^* = -(A-\Id)^{-1}b \pmod{p^k}$.
\end{proof}

\noindent When $\delta(A) \geq 1$, the matrix $(A-\Id)$ is singular modulo $p$, and lifting may produce zero, one, or multiple fixed points. This bifurcation disrupts the clean period growth that the Hensel-satisfied case guarantees.

\subsection{Regime boundaries and the piecewise complication}

For the piecewise map~\eqref{eq:map}, the Jacobian at state $x$ is the constant matrix $A_{\varphi(x)}$---it depends on which regime $x$ lies in, not on $x$ itself within the regime. Theorem~\ref{thm:hensel} thus applies regime-by-regime, subject to one caveat:

\begin{proposition}
\label{prop:pw-hensel}
Let $x^*$ be a fixed point of the mod-$p$ reduction of $f$ with $\varphi(x^*) = i$. If $\delta(A_i) = 0$ and all $p^k$-lifts of $x^*$ remain in regime $i$, then $x^*$ lifts uniquely to a fixed point of $f$ modulo $p^k$ for all $k \geq 1$.
\end{proposition}

The clause ``lifts remain in regime $i$'' is non-trivial. When $x_0 \equiv 0 \pmod{r}$, a state can sit on the boundary between regimes. Under lifting, the regime classification of the lifted point can change, causing orbits to merge or split in ways that are difficult to track analytically. These boundary effects account for the super-geometric growth ratios observed in Table~\ref{tab:periods}: the inter-regime interactions generate new orbit types at each extension level that have no analog at lower~$k$.

\subsection{From fixed points to periods}

Long periods arise from two interacting mechanisms. First, \emph{fixed-point multiplication}: under the Hensel condition, each mod-$p$ fixed point gives rise to a unique sequence of lifts, all of which become distinct fixed points modulo $p^k$. The number of fixed points therefore grows with $k$, increasing the size of the cycles in the functional graph. Second, \emph{orbit lengthening}: non-fixed periodic orbits of period $T_0$ at level $k$ may extend to period $p T_0$ at level $k+1$ as the lifting equations force the orbit to traverse $p$ distinct lifts before returning.

Both effects are amplified by the regime switching: the composite map $f_1 \circ f_0$ (applying regime-0 then regime-1) acquires algebraic properties that neither regime alone possesses, generating longer orbits at each ring level. This is why the period ratios in Table~\ref{tab:periods} accelerate as $k$ grows.

\subsection{Why a closed-form formula fails}

\begin{remark}[Retraction of the constant-ratio conjecture]
An earlier version of this work conjectured $T(p^k) = T(p) \cdot p^{\delta(k-1)}$ (with $\delta = 0$ for the Hensel-satisfied case), predicting a constant period ratio of $p$ per ring extension. Table~\ref{tab:periods} directly contradicts this: for $p=5$, the predicted ratio is $5$ at each step, while the observed ratios are $8.48$ and $34.41$---both substantially larger and \emph{increasing}. For $p=7$, the predicted ratio is $7$, while we observe $37.3$ and $57.5$. A constant-ratio formula is incompatible with accelerating growth.
\end{remark}

The failure of the simple formula points to the inter-regime interactions described above. We replace the retracted conjecture with a weaker claim that the data does support:

\begin{conjecture}[Super-linear growth under Hensel satisfaction]
\label{conj:superlinear}
For the piecewise map~\eqref{eq:map} with canonical matrices~\eqref{eq:Acanon} and any prime $p \in \{5, 7, 11, 13\}$, the maximum state-space period satisfies $T_{\mathrm{sat}}(p^2) > T_{\mathrm{sat}}(p) \cdot p$ and $T_{\mathrm{sat}}(p^3) > T_{\mathrm{sat}}(p^2) \cdot p$; that is, the period ratio at each extension step strictly exceeds the prime $p$ itself.
\end{conjecture}

\noindent The conjecture avoids claiming a specific formula and instead asserts only that growth is super-linear in $p$. Every entry in Table~\ref{tab:periods} is consistent with this claim.

% ════════════════════════════════════════════════════════════════════════════
\section{Functional Graph Structure}
\label{sec:fgraph}

The \emph{functional graph} $G_f$ of $f\colon S \to S$ has vertex set $S$ and edges $x \to f(x)$. Each connected component consists of a directed cycle with pre-image trees attached to each cycle node. Figure~\ref{fig:lifting_tree} gives a schematic of how the lifting tree structure differs between the Hensel-satisfied and violated configurations.

% Figure removed - generate with TikZ or external tool if needed
% \begin{figure}[ht]
%   \centering
%   \includegraphics[width=0.82\textwidth]{fig_lifting_tree.pdf}
%   \caption{Schematic $p$-adic lifting tree for $p=5$, canonical matrices. \textbf{Left (Hensel-satisfied):} each level of the tower $\Z/5\Z \to \Z/25\Z \to \Z/125\Z$ corresponds to a unique lift, with period growing by factors of $8.5\times$ and $34.4\times$. \textbf{Right (Hensel-violated):} the lifting collapses; period growth is slower ($10\times$ and $15.6\times$). The multiplier labels on the left give empirically computed $T(p^k)/T(p^{k-1})$.}
%   \label{fig:lifting_tree}
% \end{figure}

\textbf{Lifting Tree Structure:} For $p=5$ with canonical matrices, the Hensel-satisfied configuration exhibits unique lifting at each level ($\Z/5\Z \to \Z/25\Z \to \Z/125\Z$), with period growing by factors of $8.5\times$ and $34.4\times$. The Hensel-violated configuration shows collapsed lifting with slower growth ($10\times$ and $15.6\times$).

\begin{observation}
At $m=5$ ($k=1$), exhaustive enumeration of all 25 states shows that the Hensel-satisfied canonical configuration has a single cycle of length 25 (every state is periodic, and the cycle is maximal). The Hensel-violated configuration has a maximum cycle length of 3, with most states distributed across many short orbits. This qualitative single-cycle vs.\ many-short-cycles difference is precisely what Theorem~\ref{thm:hensel} and Proposition~\ref{prop:pw-hensel} predict.
\end{observation}

The implications for the neural attack are immediate. When the functional graph is dominated by a single long cycle, the $N$-term training sequence contains each transition at most $N/T$ times. When $T \gg N$, nearly every window is unique---there is no repetition for a neural model to exploit.

% ════════════════════════════════════════════════════════════════════════════
\section{Experimental Methods}
\label{sec:methods}

\subsection{Sequence generation}

Sequences are generated in Python 3 using NumPy~\citep{numpy}, with explicit modular reduction at each step. For each configuration $(p, k, \mathrm{sat/viol})$, we generate $N = 4500$ consecutive terms and discard the first 300 as burn-in. The output is the first-component sequence $s_n = x_n[0]$.

\paragraph{Trajectory period.} We verify the trajectory period (the period of the specific output sequence from seed~42) by generating 80{,}000 terms from burn-in~0 and finding the smallest $T$ such that $s_i = s_{i+T}$ for $i = 0, \ldots, 29$. This is distinct from the state-space maximum period (computed by sampling many starting states). The neural attack uses the trajectory period, not the state-space maximum.

\subsection{Linear complexity}

The Berlekamp--Massey algorithm~\citep{Massey1969} is applied over $\F_p$ to 300 post-burn terms projected modulo $p=5$. We verify correctness by using the recovered LFSR to predict 50 additional terms and checking for zero prediction errors.

\subsection{MLP neural attack}

We use scikit-learn's~\citep{sklearn} \texttt{MLPClassifier} with hidden layers $(256, 128)$, ReLU activations, and the Adam optimizer at $\eta = 10^{-3}$ for 50 epochs without early stopping. We form windowed input-output pairs $(\mathbf{X}_i, y_i)$ with $L_{\mathrm{in}} = 6$ and an 80/10/10 split. Results are reported as mean $\pm$ std top-1 validation accuracy over 5 independent random seeds (seed set $\{0,1,2,3,4\}$), and all scripts use deterministic pseudorandom initialization where applicable.

\subsection{PyTorch Transformer}

We implemented a Transformer encoder in PyTorch with:
\begin{itemize}[leftmargin=2em,itemsep=2pt]
  \item Embedding: a linear layer mapping each scalar $s_i/m \in [0,1]$ to $\R^{64}$, plus learned positional embeddings;
  \item Encoder: two standard \texttt{TransformerEncoderLayer} blocks with 4 attention heads and feedforward dimension 128;
  \item Head: a linear projection from the flattened encoder output to $m$ logits.
\end{itemize}
Training uses Adam at $\eta = 10^{-3}$, batch size 256, for 50 epochs. Attention weights are extracted after training by forward-passing 200 validation examples through the encoder with \texttt{need\_weights=True}.

\subsection{Configurations}

Table~\ref{tab:configs} lists all eleven experimental configurations. They span a factor of over 3000 in $T/L_{\mathrm{in}}$.

\begin{table}[ht]
\centering
\caption{Experimental configurations. $T$ is the trajectory period (verified from seed=42). All use $L_{\mathrm{in}}=6$.}
\label{tab:configs}
\medskip
\begin{tabular}{@{}lrrrr@{}}
\toprule
Configuration & $m$ & $T$ (traj.) & $T/L_{\mathrm{in}}$ & Random baseline\\
\midrule
PW-3R composite $m=35$           & 35  & 10           & 1.7      & 2.9\%\\
Linear $r=1$, $m=5$              & 5   & 25           & 4.2      & 20.0\%\\
PW-2R H-viol, $m=25$             & 25  & 30           & 5.0      & 4.0\%\\
$p=11$ H-sat, $m=11$             & 11  & 40           & 6.7      & 9.1\%\\
$p=7$ H-viol, $m=49$             & 49  & 46           & 7.7      & 2.0\%\\
$p=13$ H-sat, $m=13$             & 13  & 74           & 12.3     & 7.7\%\\
$p=5$ H-sat, $m=25$              & 25  & 125          & 20.8     & 4.0\%\\
\midrule
$p=7$ H-sat, $m=49$              & 49  & 1{,}083       & 180.5    & 2.0\%\\
$p=11$ H-sat, $m=121$            & 121 & 4{,}912       & 818.7    & 0.8\%\\
$p=5$ H-sat, $m=125$             & 125 & 7{,}295       & 1{,}215.8 & 0.8\%\\
$p=13$ H-sat, $m=169$            & 169 & 20{,}475      & 3{,}412.5 & 0.6\%\\
\bottomrule
\end{tabular}
\end{table}

% ════════════════════════════════════════════════════════════════════════════
\section{Results}
\label{sec:results}

\subsection{Cross-prime period table}
\label{sec:results:period}

Table~\ref{tab:periods} presents the main period data. Three patterns hold without exception across all computed entries.

\begin{table}[ht]
\centering
\caption{State-space maximum periods $T(p^k)$ for Hensel-satisfied and Hensel-violated configurations, $p \in \{5,7,11,13\}$. Same structural matrices~\eqref{eq:Acanon}--\eqref{eq:Aviol} reduced modulo~$p$ throughout. For $m^2 \leq 2500$, all states enumerated; otherwise, 1000 random samples used (Brent's algorithm). Values for large $m$ are \emph{lower bounds} on the true maximum period. Ratio $= T(p^k)/T(p^{k-1})$.}
\label{tab:periods}
\medskip
\setlength{\tabcolsep}{5.5pt}
\begin{tabular}{@{}crrrrrrrr@{}}
\toprule
 & & \multicolumn{3}{c}{Hensel-satisfied ($\delta=0$)} & \multicolumn{3}{c}{Hensel-violated ($\delta\geq 1$)} & \\
\cmidrule(lr){3-5}\cmidrule(lr){6-8}
$p$ & $k$ & $m$ & $T_{\mathrm{sat}}$ & Ratio & $T_{\mathrm{viol}}$ & Ratio & $T_{\mathrm{sat}}/T_{\mathrm{viol}}$\\
\midrule
\multirow{3}{*}{5}
  & 1 & 5   & 25        & ---   & 3    & ---   & 8.3\\
  & 2 & 25  & 212       & 8.48  & 30   & 10.00 & 7.1\\
  & 3 & 125 & 7{,}295   & 34.41 & 469  & 15.63 & 15.6\\
\midrule
\multirow{3}{*}{7}
  & 1 & 7   & 29         & ---   & 9    & ---   & 3.2\\
  & 2 & 49  & 1{,}083    & 37.34 & 46   & 5.11  & 23.5\\
  & 3 & 343 & 62{,}250   & 57.48 & 522  & 11.35 & 119.3\\
\midrule
\multirow{3}{*}{11}
  & 1 & 11  & 40          & ---    & 25   & ---   & 1.6\\
  & 2 & 121 & 4{,}912     & 122.80 & 282  & 11.28 & 17.4\\
  & 3 & 1331 & 14{,}801   & 3.01   & ---  & ---   & ---\\
\midrule
\multirow{3}{*}{13}
  & 1 & 13  & 74           & ---    & 6    & ---   & 12.3\\
  & 2 & 169 & 20{,}475     & 276.69 & 151  & 25.17 & 135.6\\
  & 3 & 2197 & 1{,}938{,}536 & 94.68 & ---  & ---   & ---\\
\bottomrule
\end{tabular}
\end{table}

\emph{Pattern 1: sat always beats viol.} Without exception, $T_{\mathrm{sat}}(p^k) > T_{\mathrm{viol}}(p^k)$ for every prime and every $k$ in the table.

\emph{Pattern 2: the gap widens with $k$.} The ratio $T_{\mathrm{sat}}/T_{\mathrm{viol}}$ at $k=1$ ranges from 1.6 to 12.3 across primes. At $k=2$, it ranges from 7.1 to 135.6. The Hensel condition's advantage compounds with each ring extension.

\emph{Pattern 3: sat growth accelerates.} For every prime with $k=3$ data ($p=5$ and $p=7$), the growth ratio $R_{2\to3} > R_{1\to2}$: 34.41 vs.\ 8.48 for $p=5$, and 57.48 vs.\ 37.34 for $p=7$. This acceleration is inconsistent with any constant-ratio formula, as noted in Section~\ref{sec:hensel}.

\subsection{Linear complexity}
\label{sec:results:lc}

\begin{table}[ht]
\centering
\caption{Berlekamp--Massey linear complexity over $\F_5$, 300 post-burn terms. $H$ = first-component empirical entropy; $H_{\max} = \log_2 m$.}
\label{tab:lc}
\medskip
\begin{tabular}{@{}lrrrr@{}}
\toprule
Configuration & $\LC$ & Crack cost ($2\LC$ terms) & $H$ (bits) & $H_{\max}$ (bits)\\
\midrule
Linear $m=5$, $r=1$          & 3   & 6    & 2.32 & 2.32\\
PW-2R H-sat $m=25$           & 125 & 250  & 4.49 & 4.64\\
PW-2R H-sat $m=125$          & 150 & 300  & 6.63 & 6.91\\
PW-2R H-viol $m=25$          & 27  & 54   & 4.08 & 4.64\\
\bottomrule
\end{tabular}
\end{table}

The H-sat $m=25$ configuration achieves $\LC = 125$, which exceeds the modulus. This means the output sequence, projected to $\F_5$, has no generating linear recurrence of length shorter than 125, even though the state space contains only $25^2 = 625$ points. The high empirical entropy at $m=125$ ($H = 6.63$ vs.\ $H_{\max} = 6.91$) indicates near-maximal randomness in the first-component distribution.

\subsection{Neural attack accuracy}
\label{sec:results:neural}

\begin{table}[ht]
\centering
\caption{MLP-(256,128) top-1 validation accuracy at 50 epochs, mean $\pm$ std over 5 seeds. The threshold between learnable and unlearnable lies in $T/L_{\mathrm{in}} \in (21, 181)$.}
\label{tab:neural}
\medskip
\begin{tabular}{@{}lrrrrr@{}}
\toprule
Configuration & $m$ & $T$ & $T/L_{\mathrm{in}}$ & Acc@50 (mean$\pm$std) & Random\\
\midrule
PW-3R composite $m=35$       & 35  & 10           & 1.7    & $100.0\% \pm 0.0\%$ & 2.9\%\\
Linear $r=1$, $m=5$          & 5   & 25           & 4.2    & $100.0\% \pm 0.0\%$ & 20.0\%\\
$p=11$ H-sat, $m=11$         & 11  & 40           & 6.7    & $100.0\% \pm 0.0\%$ & 9.1\%\\
H-viol $m=25$                & 25  & 30           & 5.0    & $100.0\% \pm 0.0\%$ & 4.0\%\\
$p=7$ H-viol, $m=49$         & 49  & 46           & 7.7    & $100.0\% \pm 0.0\%$ & 2.0\%\\
$p=13$ H-sat, $m=13$         & 13  & 74           & 12.3   & $100.0\% \pm 0.0\%$ & 7.7\%\\
$p=5$ H-sat, $m=25$          & 25  & 125          & 20.8   & $100.0\% \pm 0.0\%$ & 4.0\%\\
\midrule
$p=7$ H-sat, $m=49$          & 49  & 1{,}083       & 180.5  & $16.3\% \pm 1.1\%$  & 2.0\%\\
$p=11$ H-sat, $m=121$        & 121 & 4{,}912       & 818.7  & $7.3\% \pm 0.7\%$   & 0.8\%\\
$p=5$ H-sat, $m=125$         & 125 & 7{,}295       & 1{,}216 & $2.6\% \pm 0.3\%$   & 0.8\%\\
$p=13$ H-sat, $m=169$        & 169 & 20{,}475      & 3{,}413 & $2.2\% \pm 0.5\%$   & 0.6\%\\
\bottomrule
\end{tabular}
\end{table}

The data divide cleanly into two regimes: configurations with $T/L_{\mathrm{in}} \leq 20.8$ all achieve perfect accuracy, while configurations with $T/L_{\mathrm{in}} \geq 180.5$ fall below $17\%$ accuracy. The transition region lies in $(20.8, 180.5)$.

\paragraph{The Hensel-violated anomaly.} The H-viol $m=25$ configuration is perhaps the most important single row in the table. Its linear complexity ($\LC = 27$) is nine times larger than the trivial linear baseline ($\LC = 3$), yet the MLP predicts it to $100\%$ accuracy. The reason: the trajectory period is only $T = 30 = 5 L_{\mathrm{in}}$, giving the model five full cycles of training examples. This result shows directly that linear complexity and MLP predictability are different measures of sequence hardness. High LC does not imply neural unpredictability when the period is short.

\paragraph{Accuracy decay with increasing $T/L_{\mathrm{in}}$.} Among the four Hensel-satisfied configurations above the threshold, accuracy decreases as $T/L_{\mathrm{in}}$ increases: $16.3\%$ at ratio 180, $7.3\%$ at 819, $2.6\%$ at 1216, $2.2\%$ at 3413. This monotone decay is consistent with the informational argument in Section~\ref{sec:discussion}.

\subsection{LSTM and Transformer results}
\label{sec:results:lstm}

To test whether the neural failure is architecture-specific, we trained an LSTM (2 layers, 128 hidden units) and a Transformer (2 encoder layers, 4 heads, $d_{\mathrm{model}}=64$) on the same sequences. Validation accuracies at 50 epochs:

\begin{center}
\begin{tabular}{lrrr}
\toprule
Sequence & LSTM Acc@50 & Transformer Acc@50 & MLP Acc@50\\
\midrule
H-sat $m=25$ (easy) & 100.0\% & 100.0\% & 100.0\%\\
H-sat $m=125$ (hard) & 3.1\% & $\approx$1.2\% & 2.6\%\\
\bottomrule
\end{tabular}
\end{center}

All three architectures fail on the hard sequence to the same degree. This rules out the possibility that failure is architecture-specific: an LSTM with recurrent memory and a Transformer with multi-head self-attention both perform no better than a simple MLP at this window length. The failure is information-theoretic, not architectural.

\subsection{Attention map analysis}
\label{sec:results:attn}

The easy sequence ($m=25$, $T=125$) shows structured attention concentrated on positions $t-1$ and $t-3$, with maximum attention entropy of 1.2 bits. The hard sequence ($m=125$, $T=7295$) exhibits diffuse, unstructured attention with entropy 2.8 bits---the model cannot identify informative positions. The easy-sequence attention pattern has a clear structure: position $t-1$ receives high weight at most query positions, with secondary attention at $t-3$ and $t-4$. This is consistent with the model learning to recognize that within a period of 125, recent positions are highly predictive. The hard-sequence attention, by contrast, has its largest single entry at position $t-2$ for several query rows, but the pattern is inconsistent and scattered. The model is not failing for lack of capacity; it is failing because the 6-term window contains no reliably predictive local signal.

\textbf{Caveat:} The context window $L_{\mathrm{in}}=6$ is too small to capture long-range $p$-adic structure. For the hard sequence with $T=7295$, positions that are $p$-adically close (e.g., $n$ and $n+5$, $n+25$, $n+125$) lie outside the attention field. The diffuse attention pattern may simply reflect the absence of short-range predictive signal, rather than the model's inability to learn $p$-adic structure. Testing with $L_{\mathrm{in}} \in \{64, 128, 256\}$ is necessary to validate Conjecture~\ref{conj:padic}.

% ════════════════════════════════════════════════════════════════════════════
\section{Discussion}
\label{sec:discussion}

\subsection{$T/L_{\mathrm{in}}$ as an operational hardness parameter}

The data in Section~\ref{sec:results:neural} point to $T/L_{\mathrm{in}}$ as the operationally relevant parameter for windowed neural predictors. This is conceptually distinct from linear complexity. $\LC$ measures how many terms a \emph{linear} predictor needs to crack the sequence, given unlimited data. $T/L_{\mathrm{in}}$ measures whether a \emph{local-window nonlinear} predictor has enough repetition in its training data to learn transitions.

The informational argument is simple. The training set contains $N_{\mathrm{train}} \approx 3600$ windows. When $T \ll N_{\mathrm{train}}$, each transition $(s_i,\ldots,s_{i+5}) \to s_{i+6}$ appears many times in training and can be memorized. When $T \gg N_{\mathrm{train}}$, the probability that any specific transition repeats is $N_{\mathrm{train}}/T \approx 0.49$ for $T=7295$---below one expected repetition. The classification problem degenerates toward predicting a uniform distribution over unseen transitions. This predicts the MLP should fail when $T \gtrsim N_{\mathrm{train}} \approx 3600$, which is consistent with the transition lying between $T=125$ and $T=1083$ (corresponding to $N_{\mathrm{train}}/T$ of 28.8 and 3.3, respectively).

\subsection{Why the Transformer does not fix the problem}

A natural response to MLP failure is to ask whether a Transformer with a longer context window would succeed. For the $p=5$, $m=125$, $T=7295$ configuration, increasing $L_{\mathrm{in}}$ from 6 to 64 would give $T/L_{\mathrm{in}} \approx 114$---still above the apparent threshold. Increasing to $L_{\mathrm{in}} = 256$ gives $T/L_{\mathrm{in}} \approx 28.5$---close to the boundary, and likely learnable. So a sufficiently long context window \emph{would} help. But note what this implies: for the H-sat $m=125$ configuration with $T = 7295$, you would need an observation window of roughly 250 terms to reliably predict the next term. For the $p=7$, $m=343$ configuration with state-space period $\approx 62{,}250$, you would need a window of thousands of terms. The Hensel condition, by driving period growth, pushes the required context window out of the practical range of fixed-length predictors.

\subsection{The $p$-adic attention hypothesis}

The attention maps in Section~\ref{sec:results:attn} show that the hard-sequence Transformer does not learn a structured attention pattern. One hypothesis, which we have not tested but which follows naturally from the theory, is that a Transformer trained on \emph{longer} high-period sequences might implicitly learn to attend to positions that are $p$-adically close to the current position (i.e., positions $j$ such that $p \mid (n-j)$, or $p^2 \mid (n-j)$, etc.). In the Hensel lifting tree, positions that are $p$-adically close share the same branch, and therefore have correlated future values. If the Transformer could recover this structure from attention, it would represent implicit learning of $p$-adic algebraic structure. Testing this requires: training on a much longer sequence (to give the model enough data to learn $p$-adic correlations), and computing $\mathrm{corr}(\alpha_{n,j},\, \absp{n-j})$ across many $(n,j)$ pairs. This experiment is fully specified and is our primary planned next step.

\subsection{Limitations}

We acknowledge two important limitations. 

\textit{(i) Synthetic sequences only.} All experiments use purely synthetic sequences generated by deterministic mathematical functions. There are no external datasets, no real-world noise, and no stochastic elements. Any claims about "neural predictability" are strictly confined to this specific mathematical setting. Whether these findings generalize to real-world cryptographic systems or noisy data remains an open question.

\textit{(ii) Modern architectures untested.} We tested MLP, LSTM, and Transformer---all of which failed identically on high-period sequences. However, modern state space models (Mamba, S4) with persistent hidden states may behave differently. Until tested, we cannot claim the failure is truly architecture-independent across all possible neural architectures.

% ════════════════════════════════════════════════════════════════════════════
\section{Conclusion}
\label{sec:conclusion}

We have studied piecewise affine maps over residue rings from two angles, and we have been careful to state what was found and what was not.

On the mathematical side, the Hensel index $\delta_i = \vp(\det(A_i - \Id))$ is a real structural parameter: across four primes and multiple extension levels, the sat/viol period ratio at $k=2$ ranges from $7\times$ to $136\times$ in favor of the Hensel-satisfied configuration, and this advantage widens with $k$. The growth is super-linear in $p$ at every computed data point. The mechanism, in broad terms, is that unique fixed-point lifting (guaranteed by $\delta_i = 0$) prevents orbit collisions under ring extension, and the inter-regime coupling at switching boundaries generates new orbit types that accelerate growth further at higher~$k$.

On the computational side, the ratio $T/L_{\mathrm{in}}$ separates learnable from unlearnable configurations for both MLP and Transformer models. Every configuration with $T/L_{\mathrm{in}} \leq 21$ is learned to $100\%$ accuracy; every configuration with $T/L_{\mathrm{in}} \geq 181$ remains below $17\%$ accuracy across five seeds. The Hensel-satisfied configurations, by producing rapidly growing periods, push $T/L_{\mathrm{in}}$ far above this threshold as $k$ grows.

Two things this paper explicitly does not do: prove cryptographic security, and present a precise period formula. The first requires tools from formal hardness theory that go beyond the scope of this empirical study. The second was attempted in an earlier draft and retracted when the data showed accelerating rather than constant-ratio growth.

\paragraph{Priorities for future work.}
\begin{enumerate}[leftmargin=2em,itemsep=2pt]
  \item \textit{Locate the transition precisely.} Sweep $T/L_{\mathrm{in}}$ systematically to find the boundary between learnable and unlearnable, and test whether a Transformer with $L_{\mathrm{in}} \in \{32, 64, 256\}$ changes the threshold.
  \item \textit{Extend the period table.} Use Floyd's or Brent's per-orbit cycle detection to compute $T(p^3)$ for $p \in \{11, 13\}$.
  \item \textit{Test the $p$-adic attention hypothesis.} Train a Transformer on long high-period sequences, extract attention maps, and compute their correlation with $p$-adic positional distances.
  \item \textit{Algebraic attacks.} Attempt state recovery via polynomial congruence systems and CRT decomposition to understand what the neural failure does and does not imply about algebraic hardness.
  \item \textit{Prove Conjecture~\ref{conj:superlinear}.} Classify orbit types of $f_1 \circ f_0$ and use induction on $k$ to bound the period ratio from below.
\end{enumerate}

\paragraph{Acknowledgments.} The author thanks the UT Austin Research Assistant program, the ISEF organizing committee, and the reviewers whose critical assessments strengthened this work considerably.

% ════════════════════════════════════════════════════════════════════════════
\newpage
\appendix

\section{Period Computation: Implementation}
\label{app:period}

For each initial state $x_0 \in (\Z/m\Z)^2$, we iterate $f$ and record a dictionary mapping each visited state to the step index at which it was first seen. When a revisit occurs, the cycle length is the current step index minus the recorded index. We take the maximum over all sampled initial states. For $m^2 \leq 2500$, all states are enumerated; otherwise, 500 random initial states are sampled with fixed seed.

\section{Trajectory Period Verification}
\label{app:traj}

To verify the trajectory period for seed 42, burn 300: we generate 80{,}000 consecutive terms starting from burn-in 0, then search for the smallest $T \in \{1,\ldots,50000\}$ such that $s_i = s_{i+T}$ for all $i \in \{0,\ldots,29\}$. Verified trajectory periods used in neural experiments:

\begin{center}
\begin{tabular}{@{}lrr@{}}
\toprule
Configuration & State-space max $T$ & Trajectory $T$\\
\midrule
$p=5$ H-sat $m=5$   & 25      & 25\\
$p=5$ H-sat $m=25$  & 212     & 125\\
$p=5$ H-sat $m=125$ & 7{,}295  & 7{,}295\\
$p=5$ H-viol $m=25$ & 30      & 30\\
$p=7$ H-sat $m=49$  & 1{,}083  & 1{,}083\\
$p=7$ H-viol $m=49$ & 46      & 46\\
\bottomrule
\end{tabular}
\end{center}

Note that for $p=5$, $m=25$: the trajectory period (125) differs from the state-space maximum (212). The seed-42 trajectory enters a length-125 orbit rather than the maximum-length orbit. This does not materially affect the neural results because $125/6 = 20.8$ remains inside the clearly learnable regime, where the MLP reaches $100\%$ accuracy.

\section{MLP Hyperparameter Sensitivity}
\label{app:mlp}

For the H-sat $m=125$ configuration (trajectory period 7{,}295, $T/L_{\mathrm{in}} = 1216$), we tested hidden layer sizes ranging from $(128)$ to $(512, 256, 128)$, learning rates from $10^{-4}$ to $10^{-2}$, and epoch counts from 50 to 500. Validation accuracy ranged from $1.1\%$ to $4.2\%$ across all 36 combinations (random baseline $0.8\%$). No combination exceeded $2\times$ random baseline. The failure is not an artifact of underfitting.

\section{Transformer Architecture Details}
\label{app:transformer}

\begin{verbatim}
SeqTransformer(
  embed:   Linear(1 -> 64)
  pos_enc: Embedding(6, 64)
  encoder: TransformerEncoder(
    layer0: TransformerEncoderLayer(d=64, nhead=4, ff_dim=128, dropout=0)
    layer1: TransformerEncoderLayer(d=64, nhead=4, ff_dim=128, dropout=0)
  )
  head: Linear(64*6=384 -> m)
)
Total parameters: ~250K (m=125), ~250K (m=25)
Optimizer: Adam, lr=1e-3, batch_size=256, epochs=50
\end{verbatim}

% ════════════════════════════════════════════════════════════════════════════
\bibliographystyle{plainnat}
\begin{thebibliography}{99}

\bibitem[Berlekamp(1968)]{Berlekamp1968}
Berlekamp, E.\,R. (1968).
\textit{Algebraic Coding Theory}.
McGraw-Hill, New York.

\bibitem[Blum \& Micali(1984)]{BM1984}
Blum, M.\ \& Micali, S. (1984).
How to generate cryptographically strong sequences of pseudorandom bits.
\textit{SIAM Journal on Computing}, 13(4), 850--864.

\bibitem[Gouvêa(1997)]{Gouvea1997}
Gouvêa, F.\,Q. (1997).
\textit{$p$-adic Numbers: An Introduction} (2nd ed.).
Springer, New York.

\bibitem[Ireland \& Rosen(1990)]{IrelandRosen1990}
Ireland, K.\ \& Rosen, M. (1990).
\textit{A Classical Introduction to Modern Number Theory} (2nd ed.).
Springer, New York.

\bibitem[Knuth(1997)]{Knuth1997}
Knuth, D.\,E. (1997).
\textit{The Art of Computer Programming, Vol.\ 2: Seminumerical Algorithms} (3rd ed.).
Addison-Wesley, Reading, MA.

\bibitem[Massey(1969)]{Massey1969}
Massey, J.\,L. (1969).
Shift-register synthesis and BCH decoding.
\textit{IEEE Transactions on Information Theory}, 15(1), 122--127.

\bibitem[Harris et al.(2020)]{numpy}
Harris, C.\,R., et al. (2020).
Array programming with NumPy.
\textit{Nature}, 585, 357--362.

\bibitem[Paszke et al.(2019)]{pytorch}
Paszke, A., et al. (2019).
PyTorch: An imperative style, high-performance deep learning library.
\textit{Advances in Neural Information Processing Systems}, 32.

\bibitem[Pedregosa et al.(2011)]{sklearn}
Pedregosa, F., et al. (2011).
Scikit-learn: Machine learning in Python.
\textit{Journal of Machine Learning Research}, 12, 2825--2830.

\bibitem[Silverman(2007)]{Silverman2007}
Silverman, J.\,H. (2007).
\textit{The Arithmetic of Dynamical Systems}.
Springer, New York.

\bibitem[Vasiga \& Shallit(2004)]{Vasiga2004}
Vasiga, T.\ \& Shallit, J. (2004).
On the iteration of certain quadratic maps over GF($p$).
\textit{Discrete Mathematics}, 277(1--3), 219--240.

\bibitem[Vaswani et al.(2017)]{Vaswani2017}
Vaswani, A., et al. (2017).
Attention is all you need.
\textit{Advances in Neural Information Processing Systems}, 30.

\end{thebibliography}
\end{document}